% !PRINTED 6
\weberpage{6}
%<<CONT
wollen, ist keine stetige. Denn nehmen wir irgend einen rationalen
Bruch $\mu=m:n$, worin $m$ und $n$ keinen gemeinschaftlichen Theiler
haben und nicht beide Quadratzahlen sind, so ist $\mu$ nicht das
Quadrat eines rationalen Bruches. Denn wäre $\mu=p^2:q^2$, so würde
$mq^2=np^2$ folgen, und daraus $m=p^2$, $n=q^2$, was durch die
Voraussetzung ausgeschlossen ist. Wenn wir also einen Schnitt $(A,B)$
in der Mannigfaltigkeit $\mathfrak R$ bilden, indem wir jedes Element
$a$ von $\mathfrak R$ zu $A$ rechnen, dessen Quadrat kleiner als
$\mu$ ist, und jedes Element $b$ zu $B$, dessen Quadrat grösser als
$\mu$ ist, so ist weder in $A$ ein grösstes noch in $B$ ein kleinstes
Element enthalten, und der Schnitt $(A,B)$ wird nicht durch ein
Element in $\mathfrak R$ erzeugt.

Denn angenommen, es sei $a=p:q$ irgend ein Element in $A$, also
$p^2:q^2<m:n$ oder $np^2<mq^2$; dann nehmen wir eine natürliche Zahl
$y$ beliebig und wählen eine andere natürliche Zahl $x$ so, dass
$x>y$ und $x(mq^2-np^2)>ny(2p+1)$, woraus folgt:
\[
 x^2(mq^2-np^2)>nxy(2p+1)>n(2pxy+y^2),
\]
also auch
\[
 mq^2x^2>n(px+y)^2.
\]
Setzen wir also $a'=(px+y):qx$, so ist $a'>a$ und $a'^2<\mu$,
also $a'$ auch in $A$ enthalten; und ebenso kann man zeigen, dass es
in $B$ kein kleinstes Element giebt.

Die Mannigfaltigkeit $\mathfrak R$ kann uns aber als Ausgangspunkt
dienen, um eine stetige Menge zu construiren. Die Gesammtheit aller
Schnitte in $\mathfrak R$ ist gewiss eine Mannigfaltigkeit, die mit
$\mathfrak S$ bezeichnet sein mag. Betrachten wir zwei verschiedene
ihrer Elemente $\alpha=(A,B)$, $\alpha'=(A',B')$, so wird entweder
$A$ ein Theil von $A'$ oder $A'$ ein Theil von $A$ sein. Denn wenn
irgend

\par\vfill
\noindent\rule{0.22\linewidth}{0.4pt}

{\footnotesize\noindent
Mannigfaltigkeiten, denen die Stetigkeit zukommt, sind von
\gesperrt{Weierstrass} und G.\ \gesperrt{Cantor} gebildet.

Die Definition der Stetigkeit, wie wir sie hier nach
\gesperrt{Dedekind} zu Grunde legen, ist insofern erschöpfend, als
eine in diesem Sinne stetige Menge, wenn ihr noch die gleich zu
erörternde Eigenschaft der Messbarkeit zukommt, nicht Theil einer
reicheren stetigen Menge sein kann. Ich weiss nicht, ob diese
Eigenschaft schon irgendwo nachgewiesen ist, und hoffe, bei einer
anderen Gelegenheit darauf zurückzukommen. Ich bemerke aber, dass eine
solche Eigenschaft nur bei messbaren Mengen nachweisbar ist. Eine nur
geordnete Menge kann man immer, wie dicht sie auch sein mag, als Theil
einer noch dichteren auffassen.}
%CONT>>
